\chapter{Hodnotenie systému}
\label{chap:hodnotenie}

Táto kapitola overuje systém na dvoch úrovniach: technickej presnosti detektorov (Q1) a schopnosti systému pomáhať používateľom zlepšovať sa (Q3). Validácia hodnotiaceho modelu voči ľudskému vnímaniu (Q2) je popísaná v kapitole~\ref{chap:vyskum}.

\section{Validácia detektorov}
\label{sec:detector_validation}

\subsection{Metodika merania}

Každý detektor bol overený na sade krátkych izolovaných scenárov (30--90 sekúnd), v ktorých bola kontrolovane manipulovaná práve jedna premenná. Scenáre nahrával autor podľa písomného protokolu definujúceho podmienky záznamu (vzdialenosť od kamery, výška kamery, svetlo, neutrálne pozadie). Ostatné dimenzie boli počas každého scenára zámerné neutrálne.

Ground truth anotáciu vyhotovoval autor ručne vo formáte JSON s nasledujúcou štruktúrou:

\begin{verbatim}
{
  "video_id": "V1", "duration_s": 65.0,
  "annotations": {
    "<dim>": [{"start": 10.0, "end": 15.0}, ...]
  }
}
\end{verbatim}

Pre eventové detektory (výplňové slová) sa anotujú časové body s toleranciou $\pm$200~ms; pre segmentové detektory (pohyb, hlas) sa anotujú intervaly s toleranciou $\pm$500~ms pri zhode IoU $\geq 0{,}5$.

Metriky sú vypočítané pomocou evaluačného skriptu \texttt{evaluate.py}, ktorý porovnáva detekované intervaly s ground truth na úrovni 1-sekundových binárnych okien:

\begin{itemize}
    \item \textbf{Precision} --- podiel okien správne detekovaných z celkového počtu detekovaných okien
    \item \textbf{Recall} --- podiel okien správne detekovaných z celkového počtu anotovaných okien
    \item \textbf{F1} --- harmonický priemer presnosti a citlivosti; hodnota 0 ak GT existuje ale žiadna detekcia nespadá do anotovaného okna
    \item \textbf{IoU} --- priesečník delený zjednotením (Jaccard index) pre segmentové detektory
\end{itemize}

Výsledné metriky sú priemerom hodnôt naprieč všetkými testovacími videami daného detektora. Akceptačné prahy sú stanovené pragmaticky pre účely diplomovej práce: minimálny F1 = 0,70 pre vizuálne detektory, F1 = 0,75 pre akustické a textové detektory.

\subsection{Testovací materiál}

\begin{table}[H]
    \centering
    \begin{tabularx}{\textwidth}{|l|l|X|}
        \hline
        \textbf{Detektor} & \textbf{N scenárov} & \textbf{Popis manipulácií} \\
        \hline
        \specialrule{1.2pt}{0pt}{0pt}
        Očný kontakt & 2 & Dlhé odvrátenie pohľadu ($>$1~s); otočenie chrbtom k publiku \\
        \hline
        Pohyby rúk & 4 & Ruky pri tele; prirodzená gestikulácia; intenzívna gestikulácia; žiadna gestikulácia \\
        \hline
        Pohyb bokov & 3 & Pevný stoj; kývanie zo strany na stranu; prešľapovanie \\
        \hline
        Monotónnosť hlasu & 4 & Výrazná intonácia; priemerná konverzácia; zámerne monotónne; monotónne s občasnou intonáciou \\
        \hline
        Hlasitosť & 2 & Konštantná hlasitosť; príliš tichý hlas v časti prezentácie \\
        \hline
        Výplňové slová & 3 & Bez výplniek; občasné (1/30~s); husté (5/30~s) \\
        \hline
    \end{tabularx}
    \caption{Prehľad testovacích scenárov per detektor}
    \label{tab:detector_scenarios}
\end{table}

\subsection{Výsledky}

\begin{table}[H]
    \centering
    \begin{tabular}{|l|l|c|c|c|c|}
        \hline
        \textbf{Detektor} & \textbf{Jav} & \textbf{N} & \textbf{P} & \textbf{R} & \textbf{F1} \\
        \hline
        \specialrule{1.2pt}{0pt}{0pt}
        \texttt{volume} & Príliš tichý hlas & 2 & 0,805 & 1,000 & \textbf{0,892} \\ \hline
        \texttt{hip} & Kývanie bokov & 3 & 0,844 & 0,917 & \textbf{0,874} \\ \hline
        \texttt{pitch} & Monotónny segment & 4 & 0,567 & 0,790 & \textbf{0,819}$^\dagger$ \\ \hline
        \texttt{eye\_contact} & Pozeranie mimo publika ($>$1~s) & 2 & 0,746 & 0,778 & \textbf{0,750} \\ \hline
        \texttt{filler\_words} & Detekcia výplňového slova & 3 & 0,746 & 0,773 & \textbf{0,750} \\ \hline
        \texttt{arm\_movement} & Nedostatočná gestikulácia & 4 & 0,764 & 0,697 & \textbf{0,647} \\ \hline
        \texttt{eye\_contact} & Otočenie chrbtom & 2 & 0,800 & 0,400 & \textbf{0,400} \\ \hline
        \texttt{arm\_movement} & Nadmerná gestikulácia & 4 & 0,625 & 0,417 & \textbf{0,417} \\ \hline
    \end{tabular}
    \caption{Presnosť detektorov (priemerné metriky, bin = 1~s)}
    \label{tab:detector_f1}
\end{table}

Detektory hlasitosti (F1~=~0,892), pohybu bokov (F1~=~0,874) a monotónnosti (F1~=~0,819\,*) dosahujú najlepšie výsledky. \mbox{(*~Agregát} pitch počítaný z videí V12 a V13; video V11 bez GT udalostí prispieva iba k priemeru presnosti.) Detektor očného kontaktu --- pohľad mimo publika (F1~=~0,750) a výplňové slová (F1~=~0,750) presne dosahujú akceptačné prahy. Pod prahom zostávajú tri dimenzie: nedostatočná gestikulácia (F1~=~0,647), otočenie chrbtom (F1~=~0,400) a nadmerná gestikulácia (F1~=~0,417). Nízke výsledky gestikul­ácie reflektujú subjektívnu povahu hranice medzi prirodzeným a nedostatočným pohybom, kde empiricky nastavený prah rýchlosti zápästia (0,18) zle generalizuje na iné snímacie podmienky. Nízky recall otočenia chrbtom (0,400) je spôsobený jedným testovacím videom, kde sa prezentujúci otáčal pri pohľade na plátno v nepravidelných, krátkych intervaloch, ktoré neprekročili minimálnu dobu trvania udalosti.

\subsection{Analýza chýb}

\textbf{Výplňové slová --- false positives:} Neverbálne zvuky (kašeľ, vzdych) sú príležitostne klasifikované ako \textit{uhh}. Minimálna trvacia podmienka na voiced segment čiastočne obmedzuje tento efekt.

\textbf{Očný kontakt --- pohľad mimo publika, false negatives:} Pohľady na poznámky umiestnené pod kamerou nie sú spoľahlivo detekované --- yaw uhol zostáva v rámci audience zone, pretože hlava sa nakláňa, nie otáča. Detekcia by vyžadovala kombináciu yaw a pitch uhla s nižším pitch prahom.

\textbf{Očný kontakt --- otočenie chrbtom, nízky recall:} Krátkodobé otočenia (menej ako minimálna doba trvania udalosti) nie sú zachytené. V testovacích scenároch, kde prezentujúci opakovane nakukol na plátno, systém tieto epizódy prehliadol, čo znížilo recall na 0,400.

\textbf{Monotónnosť --- false positives:} Tiché úseky reči (unvoiced frames) môžu byť mylne označené ako monotónne, ak RMS energia klesne pod prah pred filtrovaním pomocou valid\_mask. Video V11 (priemerná konverzácia) vygenerovalo 11 sekúnd falošne pozitívnych detekcií, čo výrazne znížilo priemernú presnosť (0,567) bez dopadu na agregátne F1 (počítané len z videí s GT udalosťami).

\textbf{Gestikulácia --- variabilita a nízke F1:} Prahová hodnota rýchlosti zápästia 0,18 je optimum na ladiacej sade. Pri inej vzdialenosti kamery alebo zmene perspektívy sa absolútne hodnoty normalizovaných súradníc menia. Detekcia nadmerného pohybu rúk (F1~=~0,417) má dodatočnú príčinu: hranica medzi expresívnou a nervóznou gestikuláciou je ťažko definovateľná bez subjektívneho hodnotenia, a teda ani anotácia ground truth nie je jednoznačná. Odporúčaním pre budúci vývoj je kalibrácia prahov relatívne voči výške prezentujúceho v záznamu.

\section{Pilotná štúdia účinnosti}
\label{sec:pilot_study}

\subsection{Dizajn štúdie}

Pilotná štúdia overuje, či opakované používanie systému koreluje s merateľným zlepšením prezentačných schopností. Štúdia je realizovaná ako \textbf{single-subject longitudinal design} s $N = 5$ dobrovoľníkmi, kde každý účastník je sám sebe kontrolou.

\begin{table}[H]
    \centering
    \begin{tabular}{|l|l|}
        \hline
        \textbf{Parameter} & \textbf{Hodnota} \\
        \hline
        \specialrule{1.2pt}{0pt}{0pt}
        Počet účastníkov & 5 \\ \hline
        Počet sedení / účastník & 5 \\ \hline
        Frekvencia sedení & 1$\times$ týždenne \\ \hline
        Dĺžka prezentácie & 3--5 minút \\ \hline
        Téma & Rovnaká pre všetky sedenia jedného účastníka \\ \hline
        Trvanie štúdie & 5 týždňov \\ \hline
    \end{tabular}
    \caption{Parametre pilotnej štúdie}
    \label{tab:pilot_design}
\end{table}

Rovnaká téma bola zvolená na kontrolu \textit{content-effectu}: ak by účastník menil obsah, zlepšenie by bolo konfundované s lepšou znalosťou nového materiálu. Štúdia je explicitne pilotná a nezahŕňa kontrolnú skupinu --- záver nestanovuje príčinnú súvislosť, len naznačuje potenciál systému.

\subsection{Protokol sedenia}

Každé sedenie pozostávalo z piatich krokov:
\begin{enumerate}
    \item Účastník nahrá prezentáciu mobilnou aplikáciou (3--5 minút).
    \item Systém vygeneruje analýzu a celkové hodnotenie.
    \item Účastník si preštuduje spätnú väzbu a odporúčania (5--10 minút).
    \item Účastník vyplní 3-otázkový Likert dotazník (škála 1--5).
    \item V nasledujúcom týždni opakuje postup s rovnakou témou.
\end{enumerate}

Po záverečnom (5.) sedení každý účastník vyplnil \textbf{SUS (System Usability Scale)}\cite{sus} --- 10-otázkový štandardizovaný dotazník merajúci použiteľnosť systému (skóre 0--100, prah $\geq 68$).

\subsection{Výsledky --- objektívne metriky}

\begin{table}[H]
    \centering
    \begin{tabular}{|l|c|c|c|c|c|c|c|}
        \hline
        \textbf{Účastník} & \textbf{S1} & \textbf{S2} & \textbf{S3} & \textbf{S4} & \textbf{S5} & \textbf{$\Delta$} & \textbf{\% zmena} \\
        \hline
        \specialrule{1.2pt}{0pt}{0pt}
        P1 & 61,2 & 64,8 & 70,1 & 72,4 & 75,3 & $+14{,}1$ & $+23\,\%$ \\ \hline
        P2 & 54,7 & 57,3 & 58,9 & 63,2 & 67,8 & $+13{,}1$ & $+24\,\%$ \\ \hline
        P3 & 72,1 & 73,4 & 74,8 & 76,2 & 78,5 & $+6{,}4$ & $+9\,\%$ \\ \hline
        P4 & 48,3 & 52,7 & 55,1 & 54,8 & 58,2 & $+9{,}9$ & $+20\,\%$ \\ \hline
        P5 & 65,4 & 63,1 & 68,7 & 70,3 & 72,1 & $+6{,}7$ & $+10\,\%$ \\ \hline
        \textbf{Priemer} & 60,3 & 62,3 & 65,5 & 67,4 & 70,4 & $\mathbf{+10{,}0}$ & $\mathbf{+17\,\%}$ \\ \hline
    \end{tabular}
    \caption{Trajektória celkového skóre naprieč sedeniami (S1--S5)}
    \label{tab:pilot_scores}
\end{table}

Všetkých 5 účastníkov zaznamenalo zlepšenie celkového skóre. Priemerný nárast bol $+10{,}0$ bodov ($+17\,\%$) za 5 sedení. Najvýraznejšie zlepšenie dosiahli účastníci P1 a P2 ($+23\,\%$, resp. $+24\,\%$), najmenšie P3 ($+9\,\%$) --- P3 začínal s vyšším vstupným skóre (72,1), kde priestor na zlepšenie bol menší.

Zmeny po jednotlivých dimenziách ukazujú, v ktorých oblastiach bol pokrok najpriamejší:

\begin{table}[H]
    \centering
    \begin{tabular}{|l|c|c|c|c|}
        \hline
        \textbf{Účastník} & \textbf{$\Delta$ hlas} & \textbf{$\Delta$ plynulosť} & \textbf{$\Delta$ oč. kontakt} & \textbf{$\Delta$ telo} \\
        \hline
        \specialrule{1.2pt}{0pt}{0pt}
        P1 & $+18{,}2$ & $+8{,}4$ & $+12{,}7$ & $+4{,}1$ \\ \hline
        P2 & $+12{,}4$ & $+15{,}1$ & $+9{,}3$ & $+2{,}8$ \\ \hline
        P3 & $+5{,}2$ & $+4{,}7$ & $+6{,}1$ & $+8{,}3$ \\ \hline
        P4 & $+11{,}7$ & $+6{,}3$ & $+14{,}2$ & $+1{,}5$ \\ \hline
        P5 & $+7{,}3$ & $+9{,}8$ & $+5{,}4$ & $+3{,}2$ \\ \hline
        \textbf{Priemer} & $\mathbf{+11{,}0}$ & $\mathbf{+8{,}9}$ & $\mathbf{+9{,}5}$ & $\mathbf{+4{,}0}$ \\ \hline
    \end{tabular}
    \caption{Priemerná zmena skóre po dimenziách (S5 $-$ S1)}
    \label{tab:pilot_dimensions}
\end{table}

Najvýraznejšie zlepšenie nastalo v dimenzii hlas ($\Delta = +11{,}0$) a očný kontakt ($+9{,}5$). Dimenzía telo sa zlepšovala najmenej ($+4{,}0$), čo zodpovedá jej nízkej váhe (9\,\%) a subjektívnejšiemu charakteru hodnotenia.

\subsection{Výsledky --- subjektívne hodnotenie}

\begin{table}[H]
    \centering
    \begin{tabular}{|p{6cm}|c|c|c|c|c|}
        \hline
        \textbf{Otázka} & \textbf{S1} & \textbf{S2} & \textbf{S3} & \textbf{S4} & \textbf{S5} \\
        \hline
        \specialrule{1.2pt}{0pt}{0pt}
        Spätná väzba mi pomohla identifikovať slabiny & 3,8 & 4,0 & 4,2 & 4,3 & 4,4 \\ \hline
        Cítim, že sa moje zručnosti zlepšujú & 3,4 & 3,6 & 3,9 & 4,1 & 4,3 \\ \hline
        Spätná väzba bola zrozumiteľná a konkrétna & 4,0 & 4,1 & 4,2 & 4,3 & 4,3 \\ \hline
    \end{tabular}
    \caption{Priemerné Likert hodnotenia (1--5) per sedenie}
    \label{tab:pilot_likert}
\end{table}

Všetky tri otázky vykazujú rastúci trend. Subjektívny pocit zlepšenia (Q2) rastie najrýchlejšie --- účastníci vnímali progres výraznejšie v neskorších sedeniach, keď spätná väzba začala byť konkrétnejšia a pokrok bol merateľný.

\begin{table}[H]
    \centering
    \begin{tabular}{|l|c|c|}
        \hline
        \textbf{Účastník} & \textbf{SUS skóre} & \textbf{Klasifikácia} \\
        \hline
        \specialrule{1.2pt}{0pt}{0pt}
        P1 & 82,5 & Excellent \\ \hline
        P2 & 77,5 & Good \\ \hline
        P3 & 75,0 & Good \\ \hline
        P4 & 80,0 & Excellent \\ \hline
        P5 & 72,5 & Good \\ \hline
        \textbf{Priemer} & \textbf{77,5} & \textbf{Good} \\ \hline
    \end{tabular}
    \caption{SUS skóre --- záverečné sedenie}
    \label{tab:pilot_sus}
\end{table}

Priemerné SUS skóre 77,5 je nad akceptačným prahom benchmarku (68)\cite{sus} a zodpovedá klasifikácii \textit{Good}. Žiadny účastník neskóroval pod prahom 68, čo potvrdzuje, že systém je použiteľný pre zamýšľanú cieľovú skupinu.

\subsection{Záver štúdie}

5 z 5 účastníkov zaznamenalo zlepšenie celkového skóre za 5 týždňov. Priemerný nárast celkového skóre bol $+10{,}0$ bodov ($+17\,\%$). Subjektívny pocit kompetencie vzrástol z priemeru 3,4 na 4,3 (škála 1--5). SUS = 77,5 potvrdzuje prijateľnú použiteľnosť systému.

Výsledky pilotnej štúdie naznačujú, že systém má potenciál slúžiť ako efektívny tréningový nástroj. Zároveň je nevyhnutné uznať jej obmedzenia: pri $N = 5$ bez kontrolnej skupiny nie je možné oddeliť efekt systému od prirodzeného zlepšenia opakovaným prezentovaním rovnakej témy. Záver netvrdí príčinnú súvislosť, ale naznačuje smer pre ďalší výskum s väčšou vzorkou a kontrolnou skupinou.
